\section{Introduction}\label{sec:introduction}

Tomographic reconstruction asks what can be recovered from projections rather
than from direct cell data.  In grid tilings the feasible objects are usually
required to be disjoint exact tilings, possibly with typed projections or a
fixed tile shape.  This literature already contains reconstruction
algorithms, hardness results, exact domino models, and heterogeneous or
polyatomic formulations
\cite{HermanKuba1999,ChrobakEtAl2003,ChrobakDurr2001,DurrEtAl2000,
FrosiniSimi2004,FrosiniSimi2005,GritzmannLangfeld2020,
ChrobakEtAl2012}.  We ask a different, deliberately narrow question.
Suppose that one named copy of each selected piece is placed inside a target
and only the \emph{aggregate} row and column margins are required to match.
Overlaps and holes are allowed at this relaxed stage.  What does that margin
condition say about exact covers, and how hard is it to recognize the
possible failure modes?

The distinction matters because a correct aggregate margin vector need not
certify exact coverage cell by cell.  We separate three predicates:

\begin{enumerate}
  \item \emph{support}: a nonempty margin fiber contains an exact cover;
  \item \emph{purity}: every point of the margin fiber is an exact cover;
  \item \emph{repair}: every fiber component, under replacement of exactly
        two named placements, contains an exact cover.
\end{enumerate}

Purity implies repair, and repair implies support, but neither implication
reverses.  The strictness is witnessed by complete finite certificates for
three source-locked named libraries.  These witnesses organize the problem;
the logical implications themselves are elementary, and general
reconfiguration from infeasible states is prior art
\cite{FellowsEtAl2023}.

Our main uniform result uses the \emph{largest-piece residual area}
\[
  q=|T|-|L|,
\]
where \(L\) is a largest selected piece.  Once a legal pose of \(L\) is
fixed, the residual row and column margins have total mass \(q\).
Consequently every residual placement that can occur in the margin fiber is
confined to at most \(q\) active rows, at most \(q\) active columns, and
hence at most \(q^2\) active cells.  A translation bound then gives at most
\(h(q+1)\) poses for the largest piece and yields a
\(2^{O(q\log q)}\) bound on the complete margin fiber for fixed orientation
bound \(h\).  Enumerating that bounded fiber decides support, purity, and
exact-two repair.

\subsection{Contributions and scope}

The paper makes the following precise statements.

\begin{itemize}
  \item We define the named one-copy aggregate-margin fiber and its exact-cover
        sublocus, together with support, purity, and exact-two repair.
  \item We prove the active-carrier lemma and fixed-parameter tractability of
        all three predicates in \((q,h)\) for finite cell-list pieces.
  \item We give an exact participation-type compression for existential
        decision and compressed classification in a finite named library,
        under an explicit residual-area threshold.  Literal listing remains
        output-sensitive.
  \item We record a simple overlap energy whose strict descent is sufficient
        for repair, without claiming necessity.
  \item We use certified C32, C46, and P49 libraries to separate the three
        exactness levels.
  \item In one frozen P21 configuration we distinguish connectivity of a
        selected nonnegative fiber, integral relation-lattice generation, and
        all-right-hand-side Markov connectivity.
\end{itemize}

The only uniform mathematical headline is the active-carrier/FPT theorem.
We do not claim to introduce heterogeneous tiling tomography, pair
replacement, infeasible-state repair, or algebraic tiling connectivity.
Locked pair phenomena and binomial move connectivity for exact tilings have
already been studied \cite{TuckerFoltz2023,GrossYamzon2021}.  The algebraic
language in the P21 section is classical
\cite{DiaconisSturmfels1998,AokiTakemuraYoshida2008,
CharalambousThomaVladoiu2016}.  No priority or ``first FPT'' claim is made.

\subsection{Organization}

\Cref{sec:model} fixes the finite input contract and the exact-two graph.
\Cref{sec:exactness,sec:libraries} establish the hierarchy and its finite
witnesses.  \Cref{sec:localization,sec:fpt} prove localization and the FPT
theorem.  \Cref{sec:types-energy} treats implicit libraries and the overlap
potential.  \Cref{sec:p21} gives the bounded algebraic case study.
\Cref{sec:limits} states the stop boundary.  The appendices record the
evidence and reproduction contracts.
